\documentclass[11pt]{article}
\usepackage[T1]{fontenc}
\usepackage[margin=0.82in]{geometry}
\usepackage{amsmath,amssymb,amsthm}
\usepackage{graphicx}
\usepackage{xcolor}
\usepackage[hidelinks]{hyperref}
\usepackage{microtype}

\newtheorem*{theorem}{Theorem}
\newcommand{\F}{\mathbb F}
\newcommand{\A}{\mathrm A}
\newcommand{\B}{\mathrm B}
\newcommand{\Bl}{\mathrm B_{\lambda}}
\newcommand{\de}{\delta_e}
\newcommand{\zs}{\operatorname{zs}}

\title{No singularity phase transition for Haagerup functions on free groups}
\author{Full counterexample packet for arXiv:1705.05457}
\date{11 August 2026}

\begin{document}
\maketitle

\paragraph{Status and answer.}
\textbf{Both advertised questions have negative answers.}
Let $k\geq2$, let $\F_k$ have its standard word length, and put
$f_r(x)=r^{|x|}$ for $0<r<1$.  Then
\[
        \de\ll f_r.
\]
Consequently $f_r$ is not singular to $\A(\F_k)$ and is not singular to
$\Bl(\F_k)$ for \emph{any} $0<r<1$.  In particular this holds at
$r=(2k-1)^{-1/2}$ and throughout the supercritical interval.  Thus the
suggested singularity phase transition does not occur.

\section*{The source questions}
The final item of arXiv:1705.05457 (source PDF page~42) asks whether the
critical Haagerup function is singular to the Fourier algebra and whether
$f_r$ is singular to the reduced Fourier--Stieltjes algebra above the
critical value.  The exact source paragraph is reproduced here from the
archived PDF:

\begin{center}
\fbox{\includegraphics[width=0.94\textwidth]{source_question_crop.png}}
\end{center}

The source defines $\varphi\ll\psi$ by
$\zs(\varphi)\leq\zs(\psi)$ and $\varphi\perp\psi$ by
$\zs(\varphi)\zs(\psi)=0$ (Definition~4.7).  For a discrete group, it calls
$f$ singular when $f\perp\de$ (Definition~4.20); $\A(G)$ is the closed
translation-invariant subspace generated by $\de$.  Proposition~4.10 says
that a functional is absolutely continuous with respect to another exactly
when it is a norm limit of linear combinations of its two-sided translates.

\section*{The translate identity}
Let $S$ be the standard symmetric generating set of $\F_k$, so
$|S|=2k$, and write $q=2k-1$.  For left translation
$(L_s f)(x)=f(s^{-1}x)$, the following identity holds pointwise on $\F_k$:
\begin{equation}\label{eq:key}
 \sum_{s\in S}L_s f_r
 = (r^{-1}+qr)f_r-\frac{1-r^2}{r}\de .
\end{equation}
Equivalently,
\begin{equation}\label{eq:delta}
 \de=\frac{r}{1-r^2}
 \left((r^{-1}+qr)f_r-\sum_{s\in S}L_s f_r\right).
\end{equation}

\begin{theorem}
For every $k\geq2$ and $0<r<1$, one has $\de\ll f_r$.  Therefore
$f_r\not\perp\A(\F_k)$ and $f_r\not\perp\Bl(\F_k)$.
\end{theorem}

\paragraph{Proof intuition.}
Away from the identity, multiplying a reduced word on the left by the
$2k$ generators has a rigid pattern: one choice cancels its first letter,
while the other $2k-1$ choices add a letter.  Thus the adjacency sum of the
radial function $r^{|x|}$ is an eigenfunction relation everywhere except at
the root of the Cayley tree.  Its defect at the root is a nonzero multiple
of $\de$.  Solving for that defect puts $\de$ in the finite translate span
of $f_r$, which is exactly what the source's absolute-continuity criterion
requires.

\begin{proof}
Fix $x\neq e$ with $|x|=n$.  Exactly one $s\in S$ cancels the first letter
of $x$ in $s^{-1}x$, giving length $n-1$; the other $q$ choices give length
$n+1$.  Hence
\[
 \sum_{s\in S}f_r(s^{-1}x)=r^{n-1}+qr^{n+1}
 =(r^{-1}+qr)f_r(x).
\]
At $x=e$, the left side is $(q+1)r$, whereas
$(r^{-1}+qr)f_r(e)=r^{-1}+qr$.  Their difference is
$r-r^{-1}=-(1-r^2)/r$, proving \eqref{eq:key} and \eqref{eq:delta}.

Equation~\eqref{eq:delta} expresses $\de$ as a \emph{finite} linear
combination of $f_r$ and its left translates.  Proposition~4.10 of the
source therefore gives $\de\ll f_r$, i.e.
$0\neq\zs(\de)\leq\zs(f_r)$.  In particular
$\zs(f_r)\zs(\de)=\zs(\de)\neq0$, so $f_r\not\perp\de$ and hence is not
singular to $\A(\F_k)$.  Finally $\de\in\A(\F_k)\subseteq\Bl(\F_k)$, so the
same nonzero central-support overlap rules out singularity to
$\Bl(\F_k)$.
\end{proof}

\section*{Consequences and scope}
The source records three membership regimes: $f_r\in\A(\F_k)$ below the
critical value, $f_r\in\Bl(\F_k)$ at the critical value, and
$f_r\notin\Bl(\F_k)$ above it.  The theorem shows that this membership
threshold is \emph{not} a singularity threshold.  Above the threshold,
$f_r$ is outside $\Bl(\F_k)$ but still has nonzero central-support overlap
with its element $\de$.

No approximation or limiting argument is involved, and the result covers
the full parameter interval in which the Haagerup functions in the source
are defined.  The identity also works for $k=1$, but that abelian case is
outside the advertised free-group questions.

\paragraph{Verification.}
The word-counting proof was independently audited at the identity and at
nonidentity words.  An exact-rational checker exhaustively tests all reduced
words of length at most six for nine pairs $(k,r)$.  Source PDF pages 18--19,
24, 33--34, and 42 were checked for the definitions, translate criterion,
membership facts, and exact questions.

\paragraph{Novelty and human review.}
Searches by the exact question, by the adjacency identity, and through later
work on radial positive-definite functions found no published resolution as
of 11 August 2026.  This is therefore recorded as a new full counterexample,
subject to expert review.  A reviewer should check the translation
convention in Proposition~4.10, the root coefficient in \eqref{eq:key}, and
the interpretation of singularity for a subspace.

\begin{thebibliography}{9}
\raggedright
\bibitem{source}
P. Ohrysko and M. Wasilewski,
\emph{Spectral theory of Fourier--Stieltjes algebras},
arXiv:1705.05457v4.  See Definitions~4.7 and 4.20,
Proposition~4.10, and final question~(10).
\end{thebibliography}

\end{document}
